\section{引言}

自主水下航行器(Autonomous Underwater Vehicle, AUV)在海洋探测、水下巡检、环境监测等领域发挥着日益重要的作用\cite{fossen2011handbook}。AUV执行轨迹跟踪任务时，海洋环境中的海流是影响控制精度的主要扰动源。准确估计海流速度并将其前馈补偿到控制律中，是提升AUV抗流控制性能的关键\cite{borhaug2007model}。

当前AUV海流估计方法可分为两类范式。第一类是基于运动学的海流观测器\cite{liang2018three,he2024three}，通过位置误差或速度误差的积分来估计惯性系海流分量。这类方法不依赖水动力模型的精度，因此具有天然的鲁棒性，但其收敛速度受限于位置误差的累积过程，通常需要10--30秒才能收敛，难以快速响应海流突变。第二类是基于扩张状态观测器(Extended State Observer, ESO)的集总扰动估计方法\cite{xie2020extended}，将海流、模型不确定性、未建模动态等统一估计为一个"总扰动"信号。这类方法响应快速，但无法区分物理海流与模型误差，其输出是补偿信号而非物理海流参数。

值得注意的是，上述两类方法在"水动力模型失配"这一实际问题上的表现截然不同。运动学观测器因不使用动力学模型而天然免疫模型失配；ESO类方法将模型误差纳入集总扰动，不区分来源。然而，当研究者试图利用动力学模型直接估计物理海流参数时——如\cite{kim2018estimating}基于在线递推最小二乘辨识的高增益观测器、\cite{borhaug2007model}基于完美模型假设的非线性Luenberger观测器——模型失配会直接污染海流估计。这构成了本文的核心动机：\textbf{能否设计一种既利用动力学模型提供的丰富信息，又对模型失配具有鲁棒性的海流观测器？}

本文的回答是肯定的。我们的出发点是一个独特的工程条件：XHY AUV拥有通过计算流体力学(CFD)雷诺平均Navier-Stokes(RANS)仿真获得的预标定水动力模型，各自由度阻力系数的拟合优度$R^2>0.99$，并经过静水池拖曳试验和自由航行试验的独立验证。这意味着我们掌握的不仅是一个"名义模型"，更是一个具有量化精度上界的"确定性先验"。

基于此，本文提出激励门控CFD先验海流观测器(Excitation-Gated CFD-Prior Current Observer, EG-UCCO)。其核心创新包括：

\begin{enumerate}
    \item \textbf{CFD预标定模型作为确定性先验}：区别于在线参数辨识\cite{kim2018estimating}，利用预标定的高精度阻力模型作为已知先验，避免了在线辨识中"参数误差vs海流效应"的可辨识性困境。
    \item \textbf{灵敏度Gramian激励门控}：基于Fisher信息矩阵的条件数设计门控机制，仅在动力学可辨识窗口内更新海流估计。这比经验性的$|r|>r_{th}$门控具有更严格的数学基础。
    \item \textbf{速度层面的预测-校正架构}：通过在速度层面（而非加速度层面）进行一步前向预测和梯度校正，避免了加速度差分带来的噪声放大的问题（信噪比分析表明加速度层面信噪比仅约0.03）。
\end{enumerate}

本文在XHY AUV仿真平台上对EG-UCCO进行了系统验证，涵盖3种海流场景（圆形机动、直线低激励、海流阶跃）、4种模型失配水平（CFD阻力系数扰动0--30\%）、5种对比方法（运动学观测器KIN、标准LESO、物理信息PIESO、调优EKF）。实验结果表明，EG-UCCO在所有12个场景-失配组合中估计精度均为最优，且在30\%模型失配下估计精度仅退化4\%，显著优于调优EKF的16\%退化。门控消融实验证实，激励门控在高传感器噪声下可提升23\%的估计精度。

本文的剩余部分组织如下：第2节建立AUV动力学模型和海流参数化形式，并量化CFD模型的不确定性界；第3节详细阐述EG-UCCO的观测器设计和理论分析；第4节给出系统性的仿真验证结果；第5节讨论方法的适用场景与局限性；第6节总结全文。
